\section{Policy Implications}

The evidence does not imply that policy should simply push workers into large
firms or favor large firms as such. The paper does not observe firm productivity
or worker-firm matches, so it cannot prove that the wage premium is a causal
return to scale. The policy relevance of the results is more precise: a large
share of Colombian employment remains in solo work and micro-firms, the middle
of the firm-size distribution is thin, and larger firms pay higher wages even
after conditioning on observed worker characteristics, formality, and
sector-year conditions. This combination is a warning sign of labor-market
segmentation and possible allocative frictions.

The first implication is that policies should avoid preserving smallness as an
objective in itself. Support for small firms is most valuable when it helps
productive firms grow, formalize, adopt better technologies, and connect to
larger markets. By contrast, policies that subsidize informality or create
thresholds that make growth costly may reinforce the same employment structure
documented in the paper: many workers in very small units, few workers in the
middle, and a large wage gap between low-scale and large-scale employment.

The second implication is that formalization should be treated as a mechanism,
not only as an outcome. The attenuation of the firm-size coefficient after
adding formality shows that the firm-size gradient is strongly connected to the
formal segment of the labor market. This supports policies that lower the fixed
costs of formal operation, simplify compliance, improve enforcement where evasion
is strategic, and make social protection less dependent on remaining outside
formal employment. At the same time, because the premium survives controlling
for formality, formalization alone is unlikely to eliminate the wage gradient.

The third implication concerns worker mobility. If large and growing firms offer
better wage opportunities, workers need credible paths into those jobs. This
points to training, job-search intermediation, certification, childcare and
transport policies, and regional mobility tools that reduce barriers between
workers in low-scale activities and firms with better matches. The descriptive
wage gaps by education and sector show why these policies matter: firm size is
not the only axis of inequality, and workers are sorted across schooling groups
and sectors before firm size enters the analysis.

Finally, the results should be linked to productivity policy. The premium is
consistent with, but does not prove, a misallocation story. A stronger test would
combine household-survey wage patterns with firm-level measures of value added,
capital, employment, and productivity. In the meantime, the evidence supports a
policy agenda focused on removing barriers to productive firm growth, improving
the quality of worker-firm matches, and helping workers move from low-scale
employment toward more productive, formal, and scalable firms. This is closely
aligned with the broader growth agenda for Colombia emphasized by
\citet{eslava2026crecimiento}: the objective is not firm size per se, but the
reorganization of production toward higher-productivity employers and better
jobs.
